\documentclass[dvisvgm,tikz,border=3pt]{standalone}
\usepackage{amsmath,amssymb}
\usepackage{pgfplots}
\usepackage{CJKutf8}
\pgfplotsset{compat=1.18}
\usetikzlibrary{arrows.meta,calc}

\definecolor{themebg}{HTML}{2e362c}
\definecolor{themefg}{HTML}{edf4e8}
\definecolor{themegray}{HTML}{9ea897}
\definecolor{themecurve}{HTML}{7eb6ff}
\definecolor{themealert}{HTML}{ff7b7b}
\definecolor{themeamber}{HTML}{f5b942}

\begin{document}
\begin{CJK*}{UTF8}{gbsn}
\pagecolor{themebg}
\begin{tikzpicture}[color=themefg, text=themefg]

% Left Subplot: Injective (Single intersection)
\begin{axis}[
    name=axleft,
    width=7.5cm,
    height=7.2cm,
    axis lines=middle,
    axis line style={color=themefg, -{Stealth}},
    tick style={color=themefg},
    tick label style={color=themefg, font=\small},
    every axis label/.style={color=themefg},
    xmin=-2.4, xmax=2.5,
    ymin=-2.4, ymax=2.6,
    xtick={-2,-1,0,1,2},
    ytick={-2,-1,0,1,2},
    clip=false,
    xlabel={$x$},
    ylabel={$y$},
    title style={at={(0.5,1.08)},anchor=south,color=themefg,font=\bfseries\small},
    title={单射：水平线至多一交点},
]
  % Curve y = x^3/2
  \addplot[themecurve, line width=2.2pt, domain=-1.65:1.65, samples=200] {0.5*x^3};
  \node[text=themecurve, fill=themebg, inner sep=0.8pt, anchor=south west] at (axis cs:1.45, 1.85) {\small $y=\frac{1}{2}x^3$};

  % Horizontal line y = 1.0 -> x = cbrt(2) ~ 1.2599
  \draw[dashed, themegray, thick] (axis cs:-2.2,1.0) -- (axis cs:2.2,1.0);
  \node[text=themegray, fill=themebg, inner sep=0.8pt, anchor=south east] at (axis cs:-1.2, 1.05) {\small $y=y_0$};

  % Intersection point (1.2599, 1.0)
  \addplot[only marks, mark=*, mark size=2.8pt, fill=themecurve, draw=themefg] coordinates {(1.25992, 1.0)};
  \draw[dotted, themegray, thick] (axis cs:1.25992, 1.0) -- (axis cs:1.25992, 0);
  \node[text=themefg, fill=themebg, inner sep=0.8pt, below] at (axis cs:1.25992, -0.05) {\small $x_0$};
  \node[text=themecurve, fill=themebg, inner sep=0.8pt, anchor=south east] at (axis cs:1.05, 1.15) {\small $(x_0, y_0)$ 唯一原像};
\end{axis}

% Left subcaption outside axis
\node[font=\small\bfseries, color=themecurve, anchor=north] at ($(axleft.south)-(0, 0.45cm)$) {结论：任取 $y_0$，原像至多一个 $\implies f$ 单射；$f:D\to f(D)$ 存在逆映射};

% Right Subplot: Non-injective (Multiple intersections)
\begin{axis}[
    at={($(axleft.east)+(2.4cm,0)$)},
    anchor=west,
    name=axright,
    width=7.5cm,
    height=7.2cm,
    axis lines=middle,
    axis line style={color=themefg, -{Stealth}},
    tick style={color=themefg},
    tick label style={color=themefg, font=\small},
    every axis label/.style={color=themefg},
    xmin=-2.4, xmax=2.5,
    ymin=-0.6, ymax=3.6,
    xtick={-2,-1,0,1,2},
    ytick={0,1,2,3},
    clip=false,
    xlabel={$x$},
    ylabel={$y$},
    title style={at={(0.5,1.08)},anchor=south,color=themefg,font=\bfseries\small},
    title={非单射：水平线存在多交点},
]
  % Curve y = x^2
  \addplot[themealert, line width=2.2pt, domain=-1.75:1.75, samples=200] {x^2};
  \node[text=themealert, fill=themebg, inner sep=0.8pt, anchor=south east] at (axis cs:1.85, 2.7) {\small $y=x^2$};

  % Horizontal line y = 1.44 -> x = +-1.2
  \draw[dashed, themegray, thick] (axis cs:-2.2,1.44) -- (axis cs:2.2,1.44);
  \node[text=themegray, fill=themebg, inner sep=0.8pt, anchor=south east] at (axis cs:-1.4, 1.48) {\small $y=y_0$};

  % Intersection points (-1.2, 1.44) and (1.2, 1.44)
  \addplot[only marks, mark=*, mark size=2.8pt, fill=themealert, draw=themefg] coordinates {(-1.2, 1.44) (1.2, 1.44)};
  
  \draw[dotted, themegray, thick] (axis cs:-1.2, 1.44) -- (axis cs:-1.2, 0);
  \draw[dotted, themegray, thick] (axis cs:1.2, 1.44) -- (axis cs:1.2, 0);
  
  \node[text=themefg, fill=themebg, inner sep=0.8pt, below] at (axis cs:-1.2, -0.05) {\small $x_1$};
  \node[text=themefg, fill=themebg, inner sep=0.8pt, below] at (axis cs:1.2, -0.05) {\small $x_2$};
  \node[text=themealert, fill=themebg, inner sep=0.8pt, above] at (axis cs:0, 2.1) {\small 两个原像：$x_1\ne x_2,\ f(x_1)=f(x_2)=y_0$};
\end{axis}

% Right subcaption outside axis
\node[font=\small\bfseries, color=themealert, anchor=north] at ($(axright.south)-(0, 0.45cm)$) {结论：存在 $y_0$ 有至少两个原像 $\implies f$ 非单射};

\end{tikzpicture}
\end{CJK*}
\end{document}
